\part*{Part VI --- The range, worked}
\addcontentsline{toc}{part}{Part VI --- The range, worked}

\bigskip
\noindent
\emph{Six observation disciplines, each a different setting of the dials of
Parts I and IV, each with an example that has been run. The purpose of the part is
to make the word ``range'' concrete: these are not variations on a theme but
logics that see different things, cost different amounts, and support different
kinds of learner. Numbers are quoted from the notes and their simulations
\cite{GameNote,ComposeNote,GradedWhere,QuorumNote,ScopeNote}; nothing has been
re-run here.}
\bigskip

\section{How to read this part}
\label{sec:howread}

Three conventions recur, and it is easier to state them once.

\paragraph{Pricing.} An observation is charged by the depth of the formula that
performs it, on the schedule $\kappa(d) = 2^{d}$ per candidate examined. The
schedule is a choice, and it was varied: the qualitative results below are robust
to it and the exact boundaries are not. Reading structure the environment has
already deposited is charged at the formula's spatial depth; driving the system
forward to find something out is charged per rendezvous, and costs more.

\paragraph{Currency.} Where a learner is embedded in an ecology, results are
reported as \emph{yield per metered unit} --- how much better the learner does per
unit of observation it buys. This makes claims about theories comparable to claims
about food, which is the point.

\paragraph{Ladders.} A \emph{ladder} is a sequence of strategies, each identified
with a formula, ordered by the depth and modal complexity of that formula, with
the exactly-correct characteristic formula $\chi$ at the top. A ladder makes the
observation discipline into an independent variable.

\section{Discipline 1 --- Spatial, crisp}
\label{sec:d1}

\emph{Dials: $\Coll = \Pfin$, $\V = \B$, structural connectives only, no
modality.} This is perception: predicates about the shape of the term, with no
lookahead.

The worked example is noughts-and-crosses \cite{GameNote}, over the complete
census of $255{,}168$ plays --- $5478$ reachable positions, $765$ up to symmetry,
$4520$ non-terminal, of which $1890$ ($41.8\%$) offer a fork.

The instructive result is not about play but about \emph{encoding}. A fork is two
threats that cannot both be blocked, and expressing it requires graded separation
at $k=0$: $\mathit{Threat} \mid_{0} \mathit{Threat}$, two threats sharing no gap.
Requiring that formula to be expressible \emph{forces the board encoding}: the
grading namespace must be indexed by squares, not by lines, because two threats
with a common gap must be excluded and only a square-indexed namespace can see the
sharing.

\begin{observation}[The logic constrains the model, not only the other way round]
\label{obs:encoding}
Which predicates one wants to be able to state is a constraint on how the world is
represented. This is the first and simplest sense in which choosing an observation
discipline is a commitment rather than a convenience.
\end{observation}

The purely spatial rung is also the best purchase on the entire ladder: at modal
depth $0$ and name-level depth $0$ it returns $+0.0082$ win per metered unit,
roughly five hundred times the rung below it.

\section{Discipline 2 --- Modal, crisp}
\label{sec:d2}

\emph{Dials: as above, plus the context-labelled modality.} This is interaction:
looking ahead by driving the system.

Getting this right required the correction mentioned in \S\ref{sec:three}. The
context $\Ktx$ must be a genuine redex context, which forces an empty square to be
a receipt offering to be claimed and an occupied one a resting send, so that a ply
is exactly one rewrite and $\langle \Ktx_{s}\rangle$ is a generated primitive
rather than a derived abbreviation. Redex positions then correspond bijectively to
empty squares, and the branching factor of $\exists s$ \emph{is} the legal move
count.

\begin{table}[h]
\centering
\small
\renewcommand{\arraystretch}{1.2}
\begin{tabular}{@{}lrrr@{}}
\toprule
rung & win rate & cost & return per unit \\
\midrule
0 (random)                    & .437 & 0     & --- \\
1 (win now, non-modal)        & .668 & 40.9  & $+.0057$ \\
2 (block, non-modal)          & .798 & 77.8  & $+.0035$ \\
3 ($\exists s.\langle\langle \Ktx_{s}\rangle\rangle \mathit{Fork}$) & .831 & 134.9 & $+.0006$ \\
4 ($\exists s \forall t$, alternating)  & .836 & 245.7 & $+.00005$ \\
5 (spatial, name-level depth 0) & .918 & 255.9 & $+.0082$ \\
$\chi$ (exact)                & .873 & 679.2 & $-.0001$ \\
\bottomrule
\end{tabular}
\caption{The noughts-and-crosses ladder \cite{GameNote}. Rungs are ordered by
formula complexity, not by performance.}
\label{tab:ladder}
\end{table}

Four results, all exact over the full census.

\paragraph{Truth is dominated.} $\chi$ wins $.873$ at cost $679.2$; rung 5 wins
$.918$ at $255.9$. $\chi$ is the unique rung that never loses --- truth buys
safety, not food --- and it invades a rung-5 population only if the endowment
exceeds $12{,}474$ times the unit price. Across the whole phase grid it is optimal
in two cells.

\paragraph{Perception beats interaction, quantitatively.} The three non-modal
rungs ($1,2,5$) deliver $+0.444$ of the total $+0.481$ gain --- $92.3\%$ of the
yield --- for $87.9$ of $255.9$ units, $34.3\%$ of the cost. The two modal rungs
deliver the remaining $+0.038$ for $168.0$ units. Rung 4, the one genuinely
alternating rung, is measurably the worst purchase on the ladder.

\begin{principle}[Read or drive]
\label{prin:readdrive}
There are two routes to any predicate: \emph{read} it, if the environment has
already computed it into legible structure, priced by formula depth; or
\emph{drive} it, priced by rendezvous. This is a claim about environments, not
about learners, and it is statable only because the framework prices the grades
separately and can therefore ask who is billed.
\end{principle}

\paragraph{Confinement, measured.} The best opening reverses with theory depth:
rung 5 prefers the centre ($.833$ against the corner's $.597$), $\chi$ prefers the
corner ($.896$ against the centre's $.548$). Deepening the theory while holding the
opening fixed takes the learner from $.833$ to $.548$. The corner-opening mutant
invades at identical cost and is exactly the distance-$2^{-1}$ move that
Confinement denies an individual.

\paragraph{A solved science is a dead ecology.} A monoculture of $\chi$, and a
monoculture of centre-opening rung 5, both draw every game, so the harvest rate is
$0.000$ and the population starves on a finite wall. The corner-opening rung-5
population sustains $.333$ and lives. Imperfection is load-bearing.

\paragraph{And the inversion.} Run the same construction on Nim \cite{ComposeNote}
and the headline reverses. In $\mathrm{Nim}(1,2,3)$ the rungs give win rates
$.5000/.5917/.7958/.9000$ at costs $0/15.2/31.4/38.5$, with returns
\emph{rising}: $+.00602$, $+.01262$, $+.01476$. Here $\chi$ --- move to nim-sum
zero, a name-level depth-$0$ formula --- is the cheapest rung on the ladder and
the best buy on it. In $\mathrm{Nim}(3,4,5)$ the returns are flat.

\begin{observation}[Affordability is a property of the world]
\label{obs:afford}
Whether truth is worth its price is not a fact about the learner. Nothing about
the learner changed between the two games. What changed was the environment's
structural regularity, and with it whether the correct theory is cheap to state.
\end{observation}

\section{Discipline 3 --- Graded, real-valued}
\label{sec:d3}

\emph{Dials: $\V = [0,1]$ with probabilistic-logic values, clause value the
product $s \cdot c$ of strength and confidence.} This is the discipline in which
a learner can be disappointed.

The world \cite{GradedWhere}: specimens are rich ($\sigma_{R} = 120$) or decoys
($\sigma_{D} = 0$); opening costs $\kappa_{\mathrm{break}} = 40$; a depth-2
structural test $\chi$ costs $\kappa_{\mathrm{sense}} = 4$; grazing yields $5$;
endowment $120$; horizon $400$. World A has $P(\text{rich}) = .40$,
$P(\chi \mid R) = .8$, $P(\chi \mid D) = .3$, so the posterior is $.64$ and
$E[\text{net} \mid \text{open}] = +36.80$. World B has $P(\text{rich}) = .15$ and
$P(\chi \mid D) = .4$, posterior $.26$, $E = -8.70$: opening on the test is a
losing proposition.

\paragraph{The bootstrap trap.} A learner with no prior evidence has confidence
zero, hence clause value zero, hence a permanently enabled receipt that never
fires and therefore never acquires evidence: $0.00$ openings and terminal wealth
$520.0$, the grazing baseline, against $3543.7$ with a single unit of prior
evidence. This is Proposition~\ref{prop:zeroabsorb} in the wild, and its remedy is
Definition~\ref{def:explore}: exploration is a disjunct.

\paragraph{Confidence buys variance, not yield.} In world A at scepticism
parameter $k=4$: the crisp forager has survival $.8501$ and terminal $4084.7$ ---
the highest terminal wealth measured and the lowest survival --- against the graded
forager's $.9485$ and $3543.7$.

\paragraph{And in a harsh world it is not variance but existence.} In world B the
crisp forager's survival is $.0029$. It cannot learn to stop, because there is
nowhere for the disappointment to go: the formula $\chi$ was true and remains
true. The graded forager's policy is the formula's \emph{value}, and the value
moves --- strength $.750 \to .250$, confidence $.333 \to .600$, product
$.250 \to .150$, probability of opening $.333 \to .231$ over five openings.

\paragraph{Amortisation, computed.} Against the true posterior $.6400$: a learner
whose hypotheses are about individuals, so that evidence cannot accumulate past
$n=2$, has mean error $.2039$; a learner whose hypotheses are about classes
reaches $n = 101.54$ and mean error $.0383$ --- better by a factor of $5.3$, for
evidence that costs nothing extra. This is the defensible version of a claim the
original note overstated: class hypotheses have an amortisation advantage, not a
monopoly. Individual hypotheses are not valueless; they outlive the crisp forager
in both worlds.

\paragraph{Scepticism has a metabolic optimum.} Under log terminal wealth --- the
right criterion for a multiplicative absorbing process --- world A peaks at
$k^{*} = 9$ ($7.8741$, against $7.744$ at $k=4$ and $7.715$ at $k=20$), and a
$50$/$50$ ensemble of the two worlds peaks at $k^{*} = 64$ ($6.4309$, against
$6.343$ at $32$ and $6.263$ at $512$). World B alone runs to the boundary, because
abstinence is right there.

\begin{principle}[Scepticism is priced by the environment]
\label{prin:scept}
A parameter that in the source logic is set by taste --- how much evidence it
takes to become confident --- acquires a selection criterion once the learner is
mortal, and the optimum rises roughly sevenfold with the riskiness of the world.
\end{principle}

\begin{remark}[Confidence needs an outside option to be visible]
If the competing branch is the complement of the same belief, confidence cancels
exactly: $sc / (sc + (1-s)c) = s$. Grading a guard buys nothing unless the learner
has something else to do. In the forager that something else is grazing.
\end{remark}

\section{Discipline 4 --- Graded, complex-valued}
\label{sec:d4}

\emph{Dials: $\V = \Cx$, clauses isometric, modality-free.} This is the coherent
regime, and its example is Shor's algorithm written in the guard language
\cite{GradedWhere}.

The construction is the join gadget of Theorem~\ref{thm:scattering}: a qubit is
which datum sits on a channel; a gate is a join over the input channel and a
two-datum ancilla channel with both ancilla patterns joined, so everything is
consumed and the outcomes coincide across input states; the clause supplies the
matrix entries. Two-qubit gates take four candidates and an arbitrary $4\times4$,
and entangle. Gadgets compose in series on fresh channels, and the term's causal
order \emph{is} the circuit's directed graph --- a layer-2 gadget is not enabled
until its input channel carries a datum, so data dependency enforces circuit order
with no scheduler.

Measured: the quantum Fourier transform assembled from gadgets matches the DFT
matrix to worst-entry error $2.7\times10^{-16}$, $1.4\times10^{-15}$,
$3.8\times10^{-15}$ at $n = 2,3,4$. Modular exponentiation preserves norm to
$1.000000000000$ on a $16$-term superposition --- deterministic communication has
one candidate of weight one, so the classical part of a quantum algorithm costs
nothing special in this semantics. For $N=15$, $a=7$ with $4+4$ qubits: eight
one-qubit and eight two-qubit gadget firings, $390$ candidate evaluations, $16$
non-zero configurations, and the counting register at $0/4/8/12$ with probability
exactly $0.250000000$ each, totalling $1.000000000000$; continued fractions give
$r=4$, and $\gcd(7^{2}-1,15) = 3$, $\gcd(7^{2}+1,15) = 5$.

\begin{observation}[Shor is an interference deficit arranged on purpose]
\label{obs:shordeficit}
Twelve of the sixteen reachable outcomes carry amplitude zero. Corollary
\ref{cor:support}'s inclusion $\mathrm{supp} \subseteq \mathrm{Reach}$ is strict by
exactly those twelve. What a quantum algorithm \emph{is}, in this semantics, is a
large deliberate strictness.
\end{observation}

The honest limits are worth repeating. The encoding is circuit-shaped: every gate
is a hand-built join gadget, and mobility, reflection and name-passing all idle.
It shows expressiveness, not a native idiom. Writing Shor is not running Shor;
classical evaluation of the denotation is exponential. And by
Proposition~\ref{prop:overshoot} the unrestricted language is $\mathsf{PP}$, so the
isometry discipline is what makes this a quantum language rather than a magic one
--- measured on a post-selection filter $\mathrm{diag}(1, 10^{-3})$ after a
Hadamard layer, where the surviving norm falls as $2^{-n/2}$ while the
renormalised probability of the all-zero outcome climbs to
$.999996$, $.999994$, $.999992$, $.999990$ for $n = 4,6,8,10$.

\section{Discipline 5 --- Semitopological}
\label{sec:d5}

\emph{Dials: $\Coll = \Open(P)$, $\V$ tropical, the coalition modality of
\S\ref{sec:semitop}.} This is the discipline of populations, and it is worked
twice: once with the observer outside the population, once inside
\cite{QuorumNote}.

\subsection{Outside: a pride reading a herd}

Model: a ring of $N = 24$ animals, neighbourhood window $2$, noise $.03$,
asynchronous majority updating. Opens are unions of contiguous arcs of length at
least $k$ --- a genuine semitopology, since two length-$k$ arcs can meet in a set
that is not open. Pricing $\kappa(d) = 2^{d}$: a heading costs $2$, deposited
confidence $4$, testing an open $+4$, one driven tick $8$.

Seven rungs, over $4000$ trials at $t=90$, $k=13$, reported as
coverage / precision / accuracy / cost: $R_{0}$ $1.0000/.4985/.4985/0$; $R_{1}$
$1.0000/.9002/.9002/2$; $R_{5}$ $.8498/.9744/.9031/4$; $R_{2}$
$1.0000/.9868/.9868/10$; $R_{3}$ $.7510/1.0000/.8755/190.2$ (return $-0.00062$);
$R_{4}$ identical at $192$; $R_{6}$ $1.0000/.9832/.9832/1344$.

\begin{observation}[The modality buys irreversibility, not accuracy]
\label{obs:irrev}
Across $6839$ coalition verdicts at five different times, precision was $1.0000$ at
every one --- including at $t=30$, where coverage is $0.37\%$. Sample majority over
the same range runs $.9160$ to $.9873$. The expensive modal rung does not predict
better; it predicts \emph{only when it is sure}, and declines otherwise.
\end{observation}

\begin{principle}[A second currency]
Yield per metered unit underprices an instrument that declines to answer. An
observer facing an irreversible commitment --- a charge, a takeoff --- wants
precision conditional on answering, not accuracy overall, and will pay a great
deal for it.
\end{principle}

\paragraph{The counter-move, and what it tells us about privacy.} Insert seams ---
zero-coupling bonds --- into the ring, from $0$ to $6$. The herd's own probability
of reaching a conclusion falls from $.7257$ to $.2331$; the pride's coverage falls
from $.7520$ to $.0115$; the pride's precision stays between $.9986$ and $1.0000$.
Illegibility is a denial of \emph{liveness}, not an induction of error --- exactly
the signature of a consensus failure seen from outside.

\begin{principle}[The legibility trade]
Safety versus liveness, and cohesion versus illegibility, are one trade-off. A
population buys privacy with the currency it buys agreement.
\end{principle}

\paragraph{Compression, and the clock.} A topen --- a set of participants
guaranteed to decide alike --- is maximal compression for an observer: inside one,
prediction accuracy is $.8683$ over $6.5$ predictions per observation; across a
seam it is $.5606$, which is chance. And the graded reading answers a question the
crisp one cannot ask: bucketing by margin quartile, the ticks-to-commitment run
$42.4/29.6/28.9/18.0$ with direction correct at
$.6080/1.0000/1.0000/1.0000$. \emph{The margin is a clock.}

\paragraph{A negative result, kept.} At $k=13$ the tropical $\otimes = \min$ over
a coalition is zeroed by a single dissenter, so the crisp and graded rungs return
identical rows. Grading pays only if the coalition size is relaxed as well: two
dials, not one. The result is kept in the note as a labelled negative.

\subsection{Inside: a swarm that decides}

Honeybee nest-site selection \cite{Seeley10,SeeleyVisscher04} is the resolutive
case: the scout is a participant evaluating a coalition formula about its own
population, and the detection \emph{is} the resolver.

\begin{proposition}[The mechanism is forced by altitude]
\label{prop:quorumforced}
Consensus-at-the-cluster is a global universal with no presented witness ---
finitely refutable but not finitely confirmable. Quorum-at-a-site is a local
$\oplus\otimes$ over one open. No bee can evaluate a global universal, so the
swarm must use the local formula.
\end{proposition}

That is Seeley and Visscher's empirical finding \cite{SeeleyVisscher04} derived
from a cost argument rather than assumed. The swarm also exhibits the split of
\S\ref{sec:weighselect} as different behaviours: dance evaluation is the clause,
piping and takeoff are the selection function, and no bee does both.

\begin{theorem}[The $2k$ theorem]
\label{thm:2k}
Two simultaneous quorums require $2k$ committed scouts. A swarm with active pool
$S < 2k$ cannot split, for any site qualities, recruitment rates, or level of
cross-inhibition.
\end{theorem}

Measured, with $k=15$: split frequency is identically $0.0000$ at
$S = 12,16,20,24,28,29$, first appears at $S = 30 = 2k$ at $0.0130$, and rises to
$0.6850$ at $S = 50$. Takeoff-with-dissent, observed occasionally by Lindauer and
by Seeley, is thereby confined to $S \ge 2k$.

\begin{principle}[Cross-inhibition changes jobs at $S = 2k$]
Below the threshold it prevents deadlock; above it, it prevents splitting.
Measured: at $S=16$, deadlock falls from $.5120$ to $.0000$ when inhibition is
switched on, with no splits in either column; at $S=40$ there is no deadlock in
either column and splitting falls from $.3860$ to $.1170$. The literature credits
stop signals with both roles and does not separate them; this is the
discriminating experiment.
\end{principle}

\paragraph{Two bottoms, and a threshold that is not a number.} A scout cannot
distinguish ``quorum not reached'' from ``I have not sampled enough''. Under
encounter sensing, false positives are impossible and all error is false negative,
and the apparent threshold scales with effective cavity area: measured $16$, $21$,
$25$, $33$ for areas $1$, $1.25$, $1.5$, $2$, giving ratios $16.0/16.8/16.7/16.5$
against a nominal $k=15$.

\paragraph{An unexpected cost, recorded honestly.} With cross-inhibition on, the
swarm picks the better site slightly \emph{less} often --- $.7235$ against $.7845$
at $q_{B} = 0.8$ --- because inhibition suppresses correct advocacy too. The
benefit of inhibition falls from $.3140$ to $.1950$ as the site disparity widens.
This was not predicted and is flagged as open.

\begin{remark}[Position statement]
None of this out-predicts the evidence-accumulator models on trajectories, and it
is not a competing dynamical model. Its contributions are a logic for what the bee
is testing, a cost argument for why it tests that rather than something else, and
structural impossibility results --- Theorem~\ref{thm:2k} --- that a statistical
model cannot state. The swarm predictions are untested proposals.
\end{remark}

\section{Discipline 6 --- Resolution-graded}
\label{sec:d6}

\emph{Dials: any of the above, with the observation depth deliberately capped.}
The question is no longer what a discipline can see but what an observer
\emph{using} it reports about a world it cannot fully resolve \cite{ScopeNote}.

Counting requires a scope, and a scope with no bound is not a computable object.
The proposal is that a \emph{namespace} supplies the computationally coherent
notion: a scope is a predicate, not a set; disjointness of scopes is structure
rather than a limitation, since $\quot{(\varphi \mid \psi)}$ with $\varphi$ and
$\psi$ disjoint parses uniquely and description length is additive; and
fractality is recovered by a fixed point on namespaces, $\mu X.\quot{((\varphi
\vee X) \mid (\psi \vee X))}$. Membership in a $\mu$-generated scope is decidable
by descent --- disjointness kills the search and the no-self-code theorem makes the
descent well-founded --- so the computational coherence of scope comes from
reflection, not from finiteness.

Copy number is then graded by the resolution $n$ at which the observer can
separate kinds, and the received integer is the flat collapse of a \emph{copy
series}. For a worked individual with four skills --- forage ($p=.95$, $d^{*}=3$),
pod ($p=.99$, $d^{*}=5$), mate ($p=.90$, $d^{*}=3$), rear ($p=.99$, $d^{*}=5$) ---
comprising $96$ learners, in a pod of $12$:

\begin{table}[h]
\centering
\small
\begin{tabular}{@{}lrr@{}}
\toprule
observer resolution & kinds reported & count inflation \\
\midrule
0 & 1 & $3.4\times$ \\
2 & --- & $2.6\times$ \\
3 & all & $1.0\times$ \\
\bottomrule
\end{tabular}
\caption{What a coarse observer reports \cite{ScopeNote}. The copy series is
$(1152, 48, 12, 1)$; its flat collapse is $1213$.}
\label{tab:observer}
\end{table}

The assembly floor for the individual is $28$ with skills sharing structure and
$43$ without --- a saving of $34.9\%$ --- and $57.1\%$ of the floor is content
rather than architecture. The most expensive skill is the least reliable one, not
the deepest.

\begin{principle}[A learner reports a world with fewer kinds in it than the world
has]
\label{prin:fewerkinds}
An observer does not fail to see the distinctions it cannot afford. It fails to
see \emph{that they are there}. The failure is invisible from inside the
discipline, and the framework can compute its size.
\end{principle}

\begin{remark}[How far the claim goes]
An earlier and stronger version of this --- ``it takes a mind to find a mind'' ---
is too strong. The observer must match the depth at which \emph{kinds} separate,
which is not the assembly index of the object. The defensible claim is about
instruments: a survey needs the second-deepest separation, since at the deepest
one the last kind comes free.
\end{remark}

\section{The measured results, collected}
\label{sec:collected}

\begin{table}[h]
\centering
\small
\renewcommand{\arraystretch}{1.25}
\begin{tabular}{@{}>{\raggedright\arraybackslash}p{5.4cm}>{\raggedright\arraybackslash}p{4.4cm}>{\raggedright\arraybackslash}p{4.2cm}@{}}
\toprule
result & figure & source \\
\midrule
Non-modal rungs dominate the ladder
  & $92.3\%$ of yield for $34.3\%$ of cost
  & \cite{GameNote} \\
The correct theory is dominated
  & $.873$ at $679.2$ vs $.918$ at $255.9$
  & \cite{GameNote} \\
Confinement, measured
  & best opening reverses; $.833 \to .548$
  & \cite{GameNote} \\
Affordability inverts with the world
  & Nim returns rising: $+.00602/+.01262/+.01476$
  & \cite{ComposeNote} \\
Bootstrap trap
  & $0.00$ openings, $520.0$ vs $3543.7$
  & \cite{GradedWhere} \\
Crisp guard cannot learn to stop
  & survival $.0029$ in the harsh world
  & \cite{GradedWhere} \\
Confidence buys variance, not yield
  & $.8501/4084.7$ vs $.9485/3543.7$
  & \cite{GradedWhere} \\
Scepticism has an interior optimum
  & $k^{*} = 9$; ensemble $k^{*} = 64$
  & \cite{GradedWhere} \\
Interleaving inflation is real
  & exactly $2$, $6$, $24$ at $n=2,3,4$
  & \cite{GradedWhere} \\
Shor as deliberate interference deficit
  & $12$ of $16$ outcomes at amplitude $0$
  & \cite{GradedWhere} \\
The modality buys irreversibility
  & precision $1.0000$ at $0.37\%$ coverage
  & \cite{QuorumNote} \\
Illegibility denies liveness, not safety
  & coverage $.7520 \to .0115$, precision held
  & \cite{QuorumNote} \\
The margin is a clock
  & $42.4/29.6/28.9/18.0$ ticks by quartile
  & \cite{QuorumNote} \\
The $2k$ theorem
  & split $0.0000$ below $S=30$; $0.0130$ at $30$
  & \cite{QuorumNote} \\
Cross-inhibition changes jobs
  & deadlock $.5120\to.0000$; split $.3860\to.1170$
  & \cite{QuorumNote} \\
A coarse observer under-reports kinds
  & $1$ kind of $4$; count inflated $3.4\times$
  & \cite{ScopeNote} \\
\bottomrule
\end{tabular}

\caption{The headline measurements, repeated from Table~\ref{tab:collected0} for
readers who came straight to this part. Every figure is produced by a simulation
shipped with its note.}
\label{tab:collected}
\end{table}
